\section{The original analytic variety and the original native scheme}

\begin{theorem}[The whole original homogeneous quotient is the existing original variety]
\label{thm:fc-original-variety-quotient-native-comparison}
\leanok
\lean{EnvelopeProperty.ProjSubvariety.quotientNativeIso, EnvelopeProperty.ProjSubvariety.quotientNativeIso_inclusion, EnvelopeProperty.ProjSubvariety.quotientNativeIso_structure, EnvelopeProperty.ProjSubvariety.quotientNativeIso_inv_structure, EnvelopeProperty.ProjSubvariety.originalNativeScheme_isIntegral, EnvelopeProperty.ProjSubvariety.originalNativeStructure_isProjectiveFactorization}
\uses{lem:fc-original-projective-variety-native-scheme}
The entire native projective scheme of the original homogeneous prime quotient
is isomorphic to the EXISTING original closed-subscheme model $X$.
This isomorphism preserves the literal original projective inclusion and
coefficient-field structure morphism, in both directions.  In particular the
existing original native scheme is integral and its original structure morphism
has an actual finite polynomial projective factorization.
\end{theorem}
\begin{proof}
The original native closed-subscheme ideal is, by its actual construction, the
kernel ideal sheaf of the original homogeneous quotient closed immersion.
Its native closed subscheme is therefore exactly that closed immersion's
scheme-theoretic image.  The canonical actual morphism from a closed immersion's
source to its scheme-theoretic image is an isomorphism.  Its defining entire
projection triangle gives the original projective inclusion equality.
The two actual coefficient maps use the same original constants, so composing
that triangle with the original coefficient projection gives the structure
morphism equality.  Composing with the inverse yields the reverse equality.
Transport the actual quotient scheme's integrality and polynomial projective
closed embedding through this whole original isomorphism.
\end{proof}

\begin{lemma}[The entire original analytic zero locus is the actual original support locus]
\label{lem:fc-original-analytic-native-support-locus}
\leanok
\lean{EnvelopeProperty.ProjSubvariety.nativeSupport_mem_original_ideal_iff, EnvelopeProperty.ProjSubvariety.originalNormalizedClosedHomeomorph}
The original analytic zero locus $X^{\mathrm{an}}$ is the entire normalized
projective valued-point locus whose actual native support belongs to the
original native ideal-sheaf support.  Every original polynomial valuation is retained.
\end{lemma}
\begin{proof}
The actual kernel ideal sheaf of the original homogeneous quotient immersion
has support equal to the original projective zero locus.  The native homogeneous
prime support of the original normalized analytic point is the kernel of its
entire original polynomial valuation.  Thus support membership says exactly
that every original defining polynomial has valuation zero, which is the original
analytic variety membership condition.  Restrict the genuine whole normalized
projective-space homeomorphism using this exact equivalence of loci.
\end{proof}

\begin{theorem}[The original analytic variety is the actual original scheme valued-point space]
\label{thm:fc-original-analytic-native-homeomorphism}
\leanok
\lean{EnvelopeProperty.ProjSubvariety.nativeAnalyticHomeomorph, EnvelopeProperty.ProjSubvariety.nativeAnalyticHomeomorph_inclusion}
\uses{lem:fc-original-analytic-native-support-locus,thm:fc-original-variety-quotient-native-comparison}
The ORIGINAL entire analytic variety $X^{\mathrm{an}}$ is canonically homeomorphic
to the actual residue-field-valued points of the EXISTING original native scheme,
lying over the trivial valuation of its original coefficient field.  Mapping
through the original native closed inclusion retains the entire original
ambient native valued point, including its residue-field valuation.
\end{theorem}
\begin{proof}
Compose the preceding whole normalized closed-locus homeomorphism with the
inverse of the actual native closed-projective valued-point homeomorphism.
The latter is obtained from the genuine residue-field closed-immersion point
lifting equivalence.  Its actual original projection identity, applied to its
inverse, gives equality of the entire mapped valued point with the original
ambient native valued point.  No comparison certificate or point lift is assumed.
\end{proof}

\begin{lemma}[Every original affine domain and polynomial value is retained]
\label{lem:fc-original-analytic-native-coordinate-values}
\leanok
\lean{EnvelopeProperty.ProjSubvariety.nativeAnalyticHomeomorph_mem_coordinateOpen, EnvelopeProperty.ProjSubvariety.nativeAnalyticHomeomorph_eval_coordinate}
\uses{thm:fc-original-analytic-native-homeomorphism}
Every original analytic coordinate chart is exactly the inverse image of its
original native projective coordinate open under the existing original inclusion.
For every original affine coordinate polynomial, its original chart seminorm
is exactly the evaluation of the corresponding entire pulled-back native
coordinate section at the original native valued point.
\end{lemma}
\begin{proof}
The whole inclusion-point identity identifies the support membership conditions
on both sides, giving the exact original domain.  The original projective-space
native comparison already preserves every whole affine coordinate polynomial
section and value.  Apply that identity to the same original ambient point.
The genuine evaluation naturality formula for the original closed immersion
identifies this evaluation with the entire pulled-back original coordinate
section on the original closed scheme.  Equality of whole valued points
preserves these evaluations, giving the stated exact original function formula.
\end{proof}
